\documentclass[a4paper,11pt]{article}
\usepackage{midterm-style}
\examinehead{企业实训课程 · 期中考试}{A 卷}
\begin{document}

\examcover{企业实训课程}{2026—2027 学年第一学期期中考试（A 卷）}{计算机组成原理、操作系统、计算机网络、数据库概论、信息系统导论（第 1—6 章）}{150 分}{180 分钟}{本卷为开卷考试，允许携带教材、笔记和纸质资料，可使用手机或电脑检索公开资料；不得相互讨论、传递答案，不得使用手机通讯或社交软件。所有答案写在答题纸上，计算题须写出主要步骤。}

% ============================================================
\parttitle{一、计算机组成原理（25 分）}
% ============================================================

\qtypename{（一）单项选择题}{5 题，每题 2 分，共 10 分}

\q 冯·诺依曼计算机体系结构的核心思想可以概括为（\quad）
\choices{A. 采用二进制、存储程序、按地址顺序执行}{B. 指令与数据分开存储、并行访问}{C. 以流水线方式重叠执行多条指令}{D. 用虚拟存储实现多道程序并发}

\q 某机字长 8 位，采用补码表示整数，则可表示的整数范围是（\quad）
\choicesline{A. $-127\sim+127$}{B. $-128\sim+127$}{C. $-128\sim+128$}{D. $0\sim255$}

\q 某计算机主频为 2.0 GHz，某程序共含 $5\times10^8$ 条指令，平均 CPI 为 1.6，则该程序的执行时间约为（\quad）
\choicesline{A. 0.2 s}{B. 0.4 s}{C. 0.8 s}{D. 1.6 s}

\q 某系统采用 Cache-主存层次，Cache 命中率为 95\%，Cache 存取周期为 2 ns，主存存取周期为 20 ns。CPU 先访问 Cache，未命中时再访问主存，且不计入调入 Cache 的时间，则平均存取时间约为（\quad）
\choicesline{A. 2.0 ns}{B. 2.9 ns}{C. 3.0 ns}{D. 3.8 ns}

\q 与程序中断方式相比，DMA 方式的主要特点是（\quad）
\choices{DMA 由硬件直接控制主存与外设之间的数据传送，传送过程不需要 CPU 逐字干预}{DMA 只能用于低速字符设备，不能用于磁盘等块设备}{DMA 每传送一个字节都必须由 CPU 执行一次中断服务程序}{DMA 不需要使用总线，因此不会与 CPU 争用总线}

\needspace{30\baselineskip}
\qtypename{（二）基础计算题}{2 题，共 8 分}

\needspace{18\baselineskip}
\q （8 分）完成下列计算，写出必要过程：
\begin{enumerate}[label=（\arabic*）, leftmargin=2.4em, itemsep=1pt]
  \item （4 分）将二进制数 $11011010.1011_2$ 分别转换为十六进制数和十进制数。
  \item （4 分）将十进制数 $-12.375$ 转换为 IEEE 754 单精度（32 位）浮点数，写出其二进制格式和十六进制表示。
\end{enumerate}
\anslines{5.5cm}

\qtypename{（三）简答题}{1 题，共 7 分}

\needspace{16\baselineskip}
\q （7 分）简述指令周期、机器周期和时钟周期的关系，并说明取指周期与执行周期的主要工作。
\anslines{6cm}

% ============================================================
\parttitle{二、操作系统（20 分）}
% ============================================================

\qtypename{（一）单项选择题}{4 题，每题 2 分，共 8 分}

\q 进程与程序的本质区别在于（\quad）
\choices{A. 进程是动态的，程序是静态的}{B. 进程占用内存，程序不占用内存}{C. 进程由 CPU 执行，程序由外设执行}{D. 进程是用户态的，程序是内核态的}

\q 银行家算法的主要目的是（\quad）
\choicesline{A. 检测死锁}{B. 避免死锁}{C. 预防死锁}{D. 解除死锁}

\q 在分页存储管理中，将逻辑地址转换为物理地址的主要依据是（\quad）
\choicesline{A. 页表}{B. 段表}{C. 文件目录}{D. 位示图}

\q 文件系统中，位示图（bitmap）通常用于（\quad）
\choicesline{A. 记录文件目录结构}{B. 管理磁盘空闲块}{C. 实现页面置换}{D. 进行地址重定位}

\needspace{30\baselineskip}
\qtypename{（二）简答题}{1 题，共 12 分}

\needspace{24\baselineskip}
\q 回答下列问题：
\begin{enumerate}[label=（\arabic*）, leftmargin=2.4em, itemsep=2pt]
  \item （6 分）说明进程的三种基本状态及其转换条件，并比较“就绪 $\to$ 运行”和“运行 $\to$ 阻塞”两种转换的主动性与触发原因。
  \item （6 分）某进程的页面访问序列为
  \[
  1,\ 2,\ 3,\ 4,\ 1,\ 2,\ 5,\ 1,\ 2,\ 3,\ 4,\ 5
  \]
  若分配给该进程 3 个页框，分别采用 FIFO 和 LRU 置换算法，计算缺页次数；再计算 4 个页框时 FIFO 的缺页次数，并据此说明什么是 Belady 异常。
\end{enumerate}
\anslines{8cm}

% ============================================================
\parttitle{三、计算机网络（35 分）}
% ============================================================

\qtypename{（一）单项选择题}{5 题，每题 2 分，共 10 分}

\q 在 TCP/IP 体系结构中，为两台主机上的应用进程提供端到端可靠字节流服务的层次是（\quad）
\choicesline{A. 网络接口层}{B. 网际层}{C. 运输层}{D. 应用层}

\q 某 CSMA/CD 局域网的数据传输速率为 10 Mb/s，总线长度为 1 km，信号传播速度为 200 m/$\mu$s。忽略中继器和网卡延迟，则该网络的最小帧长约为（\quad）
\choicesline{A. 50 bit}{B. 100 bit}{C. 200 bit}{D. 800 bit}

\q 某 TCP 连接的慢开始门限 \code{ssthresh} 初值为 8 MSS。当拥塞窗口 \code{cwnd} 增长到 8 MSS 时发生超时，则超时后 \code{ssthresh} 和 \code{cwnd} 的初值分别为（\quad）
\choicesline{A. 8 MSS，1 MSS}{B. 4 MSS，1 MSS}{C. 4 MSS，4 MSS}{D. 8 MSS，8 MSS}

\q 某 IP 数据报总长度为 4000 B（含 20 B 固定首部），需经过 MTU 为 1500 B 的链路进行分片。若分片按最大可能长度进行，则最后一个分片的片偏移字段（以 8 B 为单位）为（\quad）
\choicesline{A. 185}{B. 370}{C. 1480}{D. 2960}

\q 主机向本地域名服务器发出域名解析请求，若本地域名服务器没有相应缓存，则由其代替主机向根域名服务器继续查询。主机到本地域名服务器这种查询方式通常称为（\quad）
\choicesline{A. 递归查询}{B. 迭代查询}{C. 反向查询}{D. 广播查询}

\needspace{30\baselineskip}
\qtypename{（二）综合计算与应用题}{3 题，共 25 分}

\needspace{24\baselineskip}
\q （8 分）某 TCP 连接使用慢开始与拥塞避免算法，MSS 为 1 KB，初始 \code{cwnd} 为 1 MSS，初始 \code{ssthresh} 为 16 MSS，接收窗口足够大。请回答：
\begin{enumerate}[label=（\arabic*）, leftmargin=2.4em, itemsep=1pt]
  \item 从连接开始到 \code{cwnd} 首次达到 16 MSS，共经过几个 RTT？
  \item 若 \code{cwnd} 增长到 32 MSS 时发生超时，超时后 \code{ssthresh} 和 \code{cwnd} 分别变为多少？此后进入什么阶段？
  \item 超时恢复后，\code{cwnd} 再次增长到 16 MSS 时收到 3 个重复 ACK。按快重传/快恢复算法，\code{ssthresh} 和 \code{cwnd} 分别变为多少？收到新的确认 ACK 后 \code{cwnd} 又变为多少？
  \item 简述慢开始与拥塞避免在 \code{cwnd} 增长方式上的区别。
\end{enumerate}
\anslines{7cm}

\needspace{18\baselineskip}
\q （8 分）某 IP 数据报总长度为 4000 B（含 20 B 固定首部），DF 标志为 0。它要经过 MTU 为 1500 B 的链路。请写出分片后各分片的片偏移（以 8 B 为单位）、MF 标志和总长度，并说明分片首部中哪些字段通常保持不变、哪些字段需要修改。
\anslines{6cm}

\needspace{22\baselineskip}
\q （9 分）某单位分配到地址块 192.168.10.0/24，需要划分给 4 个部门：A 部门 100 台主机，B 部门 50 台，C 部门 20 台，D 部门 10 台。请采用 VLSM 给出一种划分方案，填写下表，并回答后面的问题。
\begin{center}\small
\begin{tabularx}{\textwidth}{@{}cXXXX@{}}
\toprule
部门 & 子网地址/掩码 & 可用主机地址范围 & 广播地址 & 可用主机数 \\
\midrule
A & & & & \\
B & & & & \\
C & & & & \\
D & & & & \\
\bottomrule
\end{tabularx}
\end{center}
\begin{enumerate}[label=（\arabic*）, leftmargin=2.4em, itemsep=1pt]
  \item 为什么不能把 192.168.10.0/24 简单等分为 4 个 /26 子网？
  \item 目的地址 192.168.10.200 属于哪个部门的子网？
\end{enumerate}
\anslines{7cm}

% ============================================================
\parttitle{四、数据库概论（35 分）}
% ============================================================

\qtypename{（一）单项选择题}{5 题，每题 2 分，共 10 分}

\q 当数据库的存储结构发生变化时，DBA 通过修改模式/内模式映像使模式保持不变，从而保证数据的（\quad）
\choicesline{A. 逻辑独立性}{B. 物理独立性}{C. 完整性}{D. 安全性}

\q 在关系代数中，查询“选修了全部课程的学生”通常需要使用（\quad）
\choicesline{A. 投影}{B. 选择}{C. 除法}{D. 自然连接}

\q 设 \code{SC} 表中课程 \code{C01} 共有 10 条选课记录，其中 2 条记录的 \code{Grade} 为空值。执行
\begin{center}
\code{SELECT COUNT(*), COUNT(Grade), AVG(Grade) FROM SC WHERE Cno='C01';}
\end{center}
后，下列叙述正确的是（\quad）
\choices{A. COUNT(*) 为 8，COUNT(Grade) 为 8}{B. COUNT(*) 为 10，COUNT(Grade) 为 8，AVG 忽略空值}{C. COUNT(*) 为 10，COUNT(Grade) 为 10，AVG 把空值按 0 计算}{D. COUNT(*) 为 8，COUNT(Grade) 为 10}

\q 下列关于范式的叙述中，正确的是（\quad）
\choices{A. 若关系模式 $R$ 属于 3NF，则 $R$ 一定属于 BCNF}{B. 若关系模式 $R$ 属于 BCNF，则 $R$ 一定属于 3NF}{C. 3NF 允许非主属性对候选码的传递函数依赖}{D. BCNF 允许主属性对候选码的部分函数依赖}

\q 事务 $T_1$ 读取某行数据后，事务 $T_2$ 修改该行并提交，$T_1$ 再次读取同一行时得到不同的值。这种现象称为（\quad）
\choicesline{A. 脏读}{B. 不可重复读}{C. 幻读}{D. 丢失修改}

\needspace{30\baselineskip}
\qtypename{（二）综合题}{3 题，共 25 分}

\needspace{30\baselineskip}
\q （10 分）设学生选课数据库中有三个基本表（加下划线的属性为主码）：
\begin{center}\small
\begin{tabular}{ll}
\toprule
关系名 & 属性 \\
\midrule
Student & \underline{Sno}, Sname, Sdept, Sage \\
Course & \underline{Cno}, Cname, Ccredit \\
SC & \underline{Sno}, \underline{Cno}, Grade \\
\bottomrule
\end{tabular}
\end{center}
请用标准 SQL 完成下列各题：
\begin{enumerate}[label=（\arabic*）, leftmargin=2.4em, itemsep=1pt]
  \item （2.5 分）查询没有选修课程 \code{C01} 的学生的学号和姓名。
  \item （2.5 分）查询平均成绩最高的学生的学号、姓名和平均成绩；若有多名学生并列最高，要求全部列出。
  \item （2.5 分）查询每门课程中成绩高于该课程平均成绩的学生的学号、课程号和成绩。
  \item （2.5 分）创建 \code{SC} 表，要求定义主码 \code{(Sno,Cno)}，定义 \code{Sno}、\code{Cno} 分别到 \code{Student}、\code{Course} 的外码，并限制 \code{Grade} 的取值在 0～100 之间。
\end{enumerate}
\anslines{9cm}

\needspace{24\baselineskip}
\q （8 分）设关系模式
\[
R(\text{学号},\text{姓名},\text{系名},\text{系主任},\text{课程号},\text{课程名},\text{成绩})
\]
其函数依赖集为
\[
F=\{\text{学号}\to\text{姓名},\ \text{学号}\to\text{系名},\ \text{系名}\to\text{系主任},\ \text{课程号}\to\text{课程名},\ (\text{学号},\text{课程号})\to\text{成绩}\}.
\]
请回答：
\begin{enumerate}[label=（\arabic*）, leftmargin=2.4em, itemsep=1pt]
  \item 写出 $R$ 的全部候选码。
  \item $R$ 最高属于第几范式？说明理由。
  \item 将 $R$ 分解为 3NF，要求分解保持函数依赖并具有无损连接性，并简要说明理由。
\end{enumerate}
\anslines{7cm}

\needspace{24\baselineskip}
\q （7 分）某关系表 \code{Orders(order\_id, customer\_id, order\_date, amount, status)} 约有 1000 万行，现有如下查询：
\begin{lstlisting}[language=SQL]
Q1: SELECT * FROM Orders
    WHERE order_date BETWEEN '2026-01-01' AND '2026-03-31'
    ORDER BY order_date;
Q2: SELECT * FROM Orders WHERE customer_id = 'C12345';
Q3: SELECT * FROM Orders WHERE status = 'UNPAID';
\end{lstlisting}
请回答：
\begin{enumerate}[label=（\arabic*）, leftmargin=2.4em, itemsep=1pt]
  \item 针对 Q1，说明在 \code{order\_date} 上建立 B+ 树索引为什么能同时改善范围选择和排序；若该索引是聚簇索引或非聚簇索引，效果有何差异？
  \item 针对 Q2，应建立什么索引？查询 \code{SELECT *} 是否一定需要回表？若只返回 \code{order\_id} 和 \code{amount}，可如何进一步优化？
  \item Q3 是否适合在 \code{status} 上建立普通 B+ 树索引？说明理由，并给出一种更合适的优化思路。
\end{enumerate}
\anslines{7cm}

% ============================================================
\parttitle{五、信息系统导论（第 1—6 章）（35 分）}
% ============================================================

\qtypename{（一）名词解释}{3 题，每题 5 分，共 15 分；要求写出概念内涵，必要时举例}

\q 信息化
\anslines{3cm}

\q 云计算
\anslines{3cm}

\q 传感器
\anslines{3cm}

\needspace{30\baselineskip}
\qtypename{（二）简答题}{1 题，共 8 分}

\needspace{16\baselineskip}
\q （8 分）简述软件工程中“软件过程”的含义，并以瀑布模型为例说明其基本阶段和主要缺点。
\anslines{5.5cm}

\needspace{36\baselineskip}
\qtypename{（三）材料分析题}{1 题，共 12 分}

\q （12 分）阅读下列材料，回答问题。
\begin{quote}\small
某农业企业建设智慧温室：在温室内部署温度、湿度、光照和 CO$_2$ 传感器，传感器数据通过 ZigBee/LoRa 网络汇聚到网关，网关再经 4G/5G 网络上传到云平台。云平台保存历史数据并运行数据分析模型，向管理人员手机推送灌溉、通风建议；同时，企业在本地保留数据库，断网时网关可缓存数据，待网络恢复后再上传。
\end{quote}
\begin{enumerate}[label=（\arabic*）, leftmargin=2.4em, itemsep=1pt]
  \item （4 分）按“感知层—网络层—平台（支撑）层—应用层”四层结构分析该系统，说明各层的主要功能。
  \item （4 分）用数据、信息、知识、智慧的层次关系说明该系统如何把传感器数据转化为管理决策，并指出数据库、计算机网络和云计算在其中的作用。
  \item （4 分）从系统可靠性和数据安全角度，各提出一条改进建议并说明理由。
\end{enumerate}
\anslines{12cm}

\vfill
\begin{center}
\textbf{—— 试卷正文结束，附页见后 ——}
\end{center}

\answerappendix{附页（一）}{如答题纸空间不足，可在此页继续作答，并注明题号。}
\answerappendix{附页（二）}{如答题纸空间不足，可在此页继续作答，并注明题号。}

\end{document}
